%; whizzy paragraph -pdf xpdf -latex ./whizzypdfptex.sh
%; whizzy-paragraph "^\\\\begin{frame}\\|\\\\emtext"
% latex beamer presentation.
% platex, latex-beamer でコンパイルすることを想定。 

%     Tokyo Debian Meeting resources
%     Copyright (C) 2012 Junichi Uekawa

%     This program is free software; you can redistribute it and/or modify
%     it under the terms of the GNU General Public License as published by
%     the Free Software Foundation; either version 2 of the License, or
%     (at your option) any later version.

%     This program is distributed in the hope that it will be useful,
%     but WITHOUT ANY WARRANTY; without even the implied warreanty of
%     MERCHANTABILITY or FITNESS FOR A PARTICULAR PURPOSE.  See the
%     GNU General Public License for more details.

%     You should have received a copy of the GNU General Public License
%     along with this program; if not, write to the Free Software
%     Foundation, Inc., 51 Franklin St, Fifth Floor, Boston, MA  02110-1301 USA

\documentclass[cjk,dvipdfmx,12pt]{beamer}
\usetheme{Tokyo}
\usepackage{monthlypresentation}

%  preview (shell-command (concat "evince " (replace-regexp-in-string "tex$" "pdf"(buffer-file-name)) "&")) 
%  presentation (shell-command (concat "xpdf -fullscreen " (replace-regexp-in-string "tex$" "pdf"(buffer-file-name)) "&"))
%  presentation (shell-command (concat "evince " (replace-regexp-in-string "tex$" "pdf"(buffer-file-name)) "&"))

%http://www.naney.org/diki/dk/hyperref.html
%日本語EUC系環境の時
\AtBeginDvi{\special{pdf:tounicode EUC-UCS2}}
%シフトJIS系環境の時
%\AtBeginDvi{\special{pdf:tounicode 90ms-RKSJ-UCS2}}

\newenvironment{commandlinesmall}%
{\VerbatimEnvironment
  \begin{Sbox}\begin{minipage}{1.0\hsize}\begin{fontsize}{8}{8} \begin{BVerbatim}}%
{\end{BVerbatim}\end{fontsize}\end{minipage}\end{Sbox}
  \setlength{\fboxsep}{8pt}
% start on a new paragraph

\vspace{6pt}% skip before
\fcolorbox{dancerdarkblue}{dancerlightblue}{\TheSbox}

\vspace{6pt}% skip after
}
%end of commandlinesmall



\title{東京エリアDebian勉強会}
\subtitle{第99回 2013年4月度}
\author{上川純一}
\date{2013年4月20日}
\logo{\includegraphics[width=8cm]{image200607/openlogo-light.eps}}

\begin{document}

\frame{\titlepage{}}

\begin{frame}{設営準備にご協力ください。}
会場設営よろしくおねがいします。
\end{frame}

\begin{frame}{Agenda}
\begin{minipage}[t]{0.45\hsize}
  \begin{itemize}
  \item 注意事項
	\begin{itemize}
	 \item 飲食禁止
	\end{itemize}
 \end{itemize}
\end{minipage} 
\begin{minipage}[t]{0.45\hsize}
 \begin{itemize}
   \item 最近あったDebian関連のイベント報告
	\begin{itemize}
        \item 第98回 東京エリアDebian勉強会
	\end{itemize}
  \item 大統一Debian勉強会2013 のご案内
  \item Debian Trivia Quiz
  \item 事前課題紹介
 \end{itemize}
\end{minipage}
\end{frame}

\section{イベント報告}
\emtext{イベント報告}
\begin{frame}{第98回 東京エリアDebian勉強会}
\end{frame}

\section{大統一Debian勉強会 2013}
\emtext{大統一Debian勉強会 2013}

\begin{frame}{大統一Debian勉強会2013}
\begin{itemize}
\item 今年も 合同Debian勉強会イベント「大統一Debian勉強会 2013」を行います。
\item 日時: 2013年6月29日（土）10:00 - 18:00 (予定)
\item 場所: 日本大学 駿河台キャンパス 駿河台校舎1号館（御茶ノ水駅から徒歩３分）
\item 今回はスポンサーを募集して行なっています。今の所、日本マイクロソフトさん、ぷらっとホームさん、トップスタジオさん、サーバーワークスさんなどが
スポンサーしていただけることが決まっています。
\item 今月28日まで発表募集中(CFP)です。詳しくは イベントWebサイト(http://gum.debian.or.jp)を参照してください。
\item 皆さんの参加、お待ちしております。
\end{itemize}

\end{frame}

\section{DWN quiz}
\emtext{DWN quiz}
\begin{frame}{Debian 常識クイズ}

Debian の常識、もちろん知ってますよね?
知らないなんて恥ずかしくて、知らないとは言えないあんなことやこんなこと、
みんなで確認してみましょう。

今回の出題範囲は\url{debian-devel-announce@lists.debian.org},
\url{debian-devel@lists.debian.org} に投稿された
内容とDebian Project Newsなどからです。

\end{frame}

\subsection{問題}
%; whizzy-master ../debianmeetingresume201304.tex
% 以上の設定をしているため、このファイルで M-x whizzytex すると、whizzytexが利用できます。
%

\santaku
{DSAのサーバを新しくホスティングしてくれるのはどれか}
{bytemark}
{grnet}
{manda}
{A}
{Debianプロジェクトのサーバ管理を担うチームがDSA、
その管理するDebianサーバのホスティングは、grnet, manda, ubcece が既存のホスティングプロバイダーで、bytemarkが今回
加わったそうです。}

\santaku
{wheezy の backports の apt-lineで適切なものはどれか}
{deb sstp://backports.debian.org/debian/ wheezy-backports backports}
{deb http://backports.debian.org/debian/ wheezy-backports main}
{deb http://ftp.debian.org/debian/ wheezy-backports main}
{C}
{backports がメインのアーカイブと統合されたようで、wheezy-backportsから
利用できるようになったようです。}

\santaku
{tech ctte 699808の結果アップロードされたのは}
{syslinux 4}
{syslinux 5}
{grub}
{A}
{新しいバージョンのsyslinuxでdebian-installerが動かないという問題があっ
てそれをダウングレードで解決したようです。}

\santaku
{あたらしくDPLになったのはだれか}
{Stephano Zacchiroli}
{Lucas Nussbaum}
{Gergely Nagy}
{B}
{Lucas が選挙の結果DPLになりました}

% \santaku
% {}
% {}
% {}
% {}
% {}
% {}


% \santaku
% {}
% {}
% {}
% {}
% {}
% {}


\section{事前課題}
\emtext{事前課題}
{\footnotesize
 %; whizzy-master ../debianmeetingresume201304.tex
% 以上の設定をしているため、このファイルで M-x whizzytex すると、
% whizzytexが利用できます

\begin{prework}{ べつやく }

\preworksection{最近Debianのアカウント・システム管理で考えていること}

ユーザアカウントの管理について特に意識したことがありません。ログイン時の
 パスワードを共通化できる仕組みがあるらしい・・・ぐらいの認識です。

\end{prework}

\begin{prework}{ koedoyoshida }

\preworksection{最近Debianのアカウント・システム管理で考えていること}

会社のサーバはLinux系で統一されているので、uid,gid,アカウント名をそれに合わせる程度です。
ldap連携も可能ですが、自分の管理しているマシンでは(社内向けのサービスを
 提供していないので)行っていません。

\preworksection{Debianでの開発環境をどう管理しているか}

stable環境にchrootでsid環境を作るorVMware環境上へ構築しています。

\preworksection{SambaをWindowsのドメインに参加させたことはありますか。}

「SambaをWindowsのドメインに参加させたこと」、「Linuxのユーザ認証をWindowsに統合したこと」はありません。
逆にsambaをDCとしてwindowsクライアントを参加させたことはあります。
\end{prework}

\begin{prework}{ ribbon@ns.ribbon.or.jp }

\preworksection{SambaをWindowsのドメインに参加させたことはありますか。}
あります。
\preworksection{参加させたことがある方、ハマったところを教えてください}
バージョンによって挙動が違うところです。
\end{prework}

\begin{prework}{ dictoss(杉本　典充) }

\preworksection{最近Debianのアカウント・システム管理で考えていること}

サーバで使っていても自分しか使用しないので、個人アカウントを作って完
 結してしまっている。

\preworksection{Debianでの開発環境をどう管理しているか}

amd64マシンを使っているので、i386環境はchrootで入れるように設定している。tarballを持ってきたアプリがビルド時にunameするものが一部あり、うまくビルドが通らないな場合を考えてKVM環境も併用している。ただ、ネットワーク関係のテストをするときはiptablesを使う場合が多いので最初からKVM環境で作って試している。
\end{prework}

\begin{prework}{ henrich }

\preworksection{最近Debianのアカウント・システム管理で考えていること}

Fedoraのようにユーザーがアカウントを取得して、ウェブからアクセスできる様々な活動ができるといいな、と思っています。例: https://admin.fedoraproject.org/updates

\preworksection{Debianでの開発環境をどう管理しているか}

どう管理、というのがちょっと意味が掴み取れません。普通にcowbuilderとpiupartsとlintian使ってるぐらいですが、これをどのように変更していけば効率の良い環境になるのかは知りたい所です。
\end{prework}

\begin{prework}{ 野首 }

\preworksection{最近Debianのアカウント・システム管理で考えていること}
Debianのアカウント管理は、自分の周りのマシンではいまだ/etc/passwdベース
 ばかりです。

\preworksection{Debianでの開発環境をどう管理しているか}
開発環境の管理としては、パッケージにsvn-buildpackageとgit-buildpackageを併用しています。特別使い分けたわけではなく、後になってgitを覚えてからgit-buildpackageを使ってみただけなので、これに特別な利点があったりはしません。むしろ設計思想が全然違うので一本化できず困っています。
あとはsid環境にlxcを使ってみたり、armやsparcのテストのためにqemuを使い始めたりしている程度です。

\end{prework}

\begin{prework}{ 石井一夫 }

\preworksection{最近Debianのアカウント・システム管理で考えていること}

セキュリティは結構深刻で、内部の者も悪さをする恐れがあり、かなり、気を使います。サーバでは、一般ユーザにsudo権限を取得できないようにし、管理者用アカウントを作ってwheel のグループに所属させ、それだけが、sudo権限を取得できるようにしています。パスワードは、3ヶ月ごとに変更。Kerberosは気になりますが、使いこなせていません。

\preworksection{Debianでの開発環境をどう管理しているか}

GCCとか、JDKとかインストールできるものは、可能な限りインストールしていますが、環境管理とまでは行っていません。これからの課題です。
\end{prework}

\begin{prework}{ 上川純一 }

\preworksection{最近Debianのアカウント・システム管理で考えていること}

個人の環境は自分と目的用途別に作っているアカウントで、自分のアカウント以
外はインタラクティブに利用する用途ではないので基本的にはパスワードを設
定せず利用しています。

昔はファイル共有でNFSを使っていたのでシステム間でUIDを一致させないといけないなどの面
倒がありましたが最近はsshfs とか rsync/ssh でファイルを共有するので片方向
の鍵認証でファイルのアクセス権が決まるモデルになっています。こっちのほう
が管理者としてはやりやすいですね。

\preworksection{Debianでの開発環境をどう管理しているか}

cowbuilderでamd64 sidの環境のみを維持管理しています。昔はいろいろなCPUアー
 キテクチャーとかOSをとかをためしていたので大量のイメージがありました。
cowbuilder だと常に最新版のバイナリを使うことになります。また、必要なソ
 フトウェアを必要なときにJust In Timeでインストールしてくれることになり
 ます。
OSイメージの構成・履歴管理などを省略してくれるのでいいです。

開発環境の仮想マシンを管理するとどういうメモリを割り当てるかとかネットワー
クの設定をどうするかだとかというのも心配することになるんですが、
cowbuilderではホストOSの環境をそのままつかっているので楽です。開発環境の
メンテナンスが目的ではなく開発環境を使うことが目的なので開発環境はメンテ
ナンスフリーなのが理想だと思います。

各種クラウドサービスを使うとOSイメージのインストールからさせられることが
多いのですがそこではじめていかに世間の人が環境の維持管理を重要なタスクだ
と思っていて、cowbuilderが楽なのかを認識した気がします。

\end{prework}

\begin{prework}{ まえだこうへい }

\preworksection{最近Debianのアカウント・システム管理で考えていること}

Debianの、というかLDAPに代わるもっとシンプルなアカウント管理の仕組みを作りたいなと思っているだけで何もやってません。

\preworksection{Debianでの開発環境をどう管理しているか}

Pythonでのツールの開発はのSidを使っていますが、Sid自体はVirtual Box上にあるので、Shutdown時に都度snapshotを取っています。

\preworksection{SambaをWindowsのドメインに参加させたことはありますか。}

某所でSamba使っていましたが、某所のADには参加できない(その某所とは別の法人組織だったので)が、アクセスは某所のPCでしかできないので、ユーザアカウントを同一にするため、某所のシステムからユーザアカウント引っこ抜いてLDIFに変換してLDAPでアカウント管理してました。(Sambaとは直接連携してない)
\end{prework}

\begin{prework}{ 吉野(yy\_{}y\_{}ja\_{}jp) }

\preworksection{最近Debianのアカウント・システム管理で考えていること}

スタンドアロンで使うとき，VMで使うときがありますがどちらも考える必要を感じてません．

\preworksection{Debianでの開発環境をどう管理しているか}

どちらの場合も特に何もしてないです．

\preworksection{SambaをWindowsのドメインに参加させたことはありますか。}
ありません．
\preworksection{Linuxのユーザ認証をWindowsに統合したことはありますか。}
ありません．

\end{prework}

\begin{prework}{ 岩松 信洋 }

\preworksection{最近Debianのアカウント・システム管理で考えていること}

特にアカウント管理に関して考えてないです。
昔ながらのユーザとグループで管理しています。

\preworksection{Debianでの開発環境をどう管理しているか}

chroot と alias, シェルスクリプトによる切り替えを使っています。

\preworksection{「SambaをWindowsのドメインに参加させたことはありますか。参加させたことがない方、もしお時間があれば実際にやってみて、ハマったところを教えてください。参加させたことがある方、ハマったところを教えてください。」、「Linuxのユーザ認証をWindowsに統合したことはありますか。ある場合はその際の方法やハマったところについて教えてください。ない場合は、もしお時間があれば実際にやってみてハマったところを教えてください」にお答えください。}

誰でもアクセスできるようなsambaサーバしか構築したことないので、特にハマったということはないです。
\end{prework}

}

\emtext{Samba で Linux の認証を Windows に統合してみたり}
\emtext{powerpc64}
\emtext{debootstrapを有効活用してみよう}
\emtext{Debian勉強会予約システム変更履歴}

\section{今後のイベント}
\emtext{今後のイベント}
\begin{frame}{今後のイベント}
\begin{itemize}
 \item 2013年5月 Debian勉強会
\end{itemize}
\end{frame}

\section{今日の宴会場所}
\emtext{今日の宴会場所}
\begin{frame}{今日の宴会場所}
未定
\end{frame}

\end{document}

;;; Local Variables: ***
;;; outline-regexp: "\\([ 	]*\\\\\\(documentstyle\\|documentclass\\|emtext\\|section\\|begin{frame}\\)\\*?[ 	]*[[{]\\|[]+\\)" ***
;;; End: ***
